\documentclass[11pt]{amsproc}

\usepackage[a4paper,margin=2.6cm]{geometry}
\usepackage{amsmath,amssymb,amsthm}
\usepackage[T1]{fontenc}
\usepackage[utf8]{inputenc}

\newcommand{\C}{\mathbb{C}}
\newcommand{\gen}[1]{\langle #1\rangle}
\newcommand{\col}{\operatorname{col}}

\begin{document}

\section*{Problema}

Pintamos cada una de las $n$ coordenadas de $\C^n$ con uno de $l$ colores. Un vector
no nulo se llama \emph{monocromático} si todas sus coordenadas no nulas son del mismo
color, y ése es entonces \emph{su} color; por ejemplo, los vectores de la base
canónica son monocromáticos.

Fijemos también una base cualquiera $c_1,\dots,c_n$ de $\C^n$ y escribamos
\[
  U_t=\gen{c_1,\dots,c_t},\qquad 0\le t\le n,
\]
de modo que $\{0\}=U_0\subset U_1\subset\cdots\subset U_n=\C^n$ y $\dim U_t=t$.

\begin{enumerate}
  \item\itshape Probar que toda base $x_1,\dots,x_n$ de $\C^n$ se puede reordenar de
        manera que los últimos $n-t$ vectores completen $U_t$:
        \[
          \C^n=U_t\oplus\gen{x_{t+1},\dots,x_n}\qquad\text{para todo } 0\le t\le n .
        \]
        Un único reordenamiento debe servir para los $n+1$ valores de $t$ a la vez.
  \item\itshape Probar que existe una base $x_1,\dots,x_n$ de $\C^n$ formada por
        vectores monocromáticos tal que, para cada $i$, sumándole a $x_i$ una
        combinación lineal de vectores $x_j$ con $j>i$ \emph{de color distinto al de
        $x_i$} se llega a un vector de $U_i$.
\end{enumerate}
 
Dicho de otro modo: en (2) cada $x_i$ alcanza el subespacio $U_i$ con la ayuda de los
vectores que vienen después, pero nunca con la ayuda de los de su propio color.

\section*{Solución}

\subsection*{Parte (1)}

Elegimos los vectores de atrás para adelante. Escribamos
\[
  T_t=\gen{x_{\sigma(t+1)},\dots,x_{\sigma(n)}}
\]
para los índices ya elegidos y
probemos, por inducción decreciente en $t$, que se los puede elegir de modo que
$\C^n=U_t\oplus T_t$.

Para $t=n$ tomamos $T_n=\{0\}$: no hay nada que hacer, pues $U_n=\C^n$.

Supongamos elegidos índices distintos $\sigma(t+1),\dots,\sigma(n)$ con
$\C^n=U_t\oplus T_t$; en particular $\dim T_t=n-t$. Como $U_{t-1}\subseteq U_t$ y
$U_t\cap T_t=\{0\}$, la suma $U_{t-1}+T_t$ es directa y
\[
  H:=U_{t-1}\oplus T_t\quad\text{tiene dimensión}\quad (t-1)+(n-t)=n-1,
\]
es decir, $H$ es un hiperplano. Los $x_1,\dots,x_n$ generan $\C^n\neq H$, así que
algún $x_i$ no pertenece a $H$. Ese índice $i$ es automáticamente uno de los que
todavía están libres: si fuera uno de los ya elegidos tendríamos
$x_i\in T_t\subseteq H$. Ponemos $\sigma(t):=i$ y obtenemos
\[
  \C^n=H\oplus\gen{x_{\sigma(t)}}=U_{t-1}\oplus\bigl(T_t\oplus\gen{x_{\sigma(t)}}\bigr)
  =U_{t-1}\oplus T_{t-1},
\]
como queríamos. Al cabo de $n$ pasos elegimos $n$ índices distintos, o sea un
reordenamiento, válido simultáneamente para todos los $t$.

\subsection*{Parte (2), paso 1}

La base canónica $e_1,\dots,e_n$ es monocromática, y esa propiedad no se pierde al
reordenar. Aplicando la primera parte a esa base obtenemos una base monocromática
$x_1,\dots,x_n$ (una permutación de la canónica) tal que
\begin{equation}\label{eq:compl}
  \C^n=U_t\oplus\gen{x_{t+1},\dots,x_n}
\end{equation}
para todo $0\le t\le n$.

\subsection*{Parte (2), paso 2}

Fijemos $i$. Por \eqref{eq:compl} con $t=i$ podemos escribir $x_i=b_i-w$ con
$b_i\in U_i$ y $w\in\gen{x_{i+1},\dots,x_n}$, o sea
\begin{equation}\label{eq:tri}
  b_i=x_i+\sum_{j>i}s_{ij}x_j\in U_i
\end{equation}
para ciertos escalares $s_{ij}\in\C$ (únicos, porque la suma en \eqref{eq:compl} es
directa). Es decir: ya sabemos llegar a $U_i$ desde $x_i$ usando vectores
posteriores. Lo único que falta es conseguir que ningún $x_j$ del color de
$x_i$ intervenga.

Los vectores $b_1,\dots,b_n$ no volverán a cambiar en lo que sigue: lo único que
modificaremos es la base $(x_i)$, y por lo tanto los coeficientes $s_{ij}$.

\subsection*{Parte (2), paso 3}

Digamos que el índice $i$ es \emph{bueno} si en \eqref{eq:tri} todos los $j$ con
$s_{ij}\neq 0$ corresponden a vectores $x_j$ de color distinto del de $x_i$. El
índice $n$ es bueno, porque $b_n=x_n$ y no hay nada que corregir.

Supongamos que los índices $i+1,\dots,n$ son buenos y arreglemos el índice $i$.
Reemplazamos $x_i$ por
\[
  x_i'=x_i+\sum_{\substack{j>i\\ \col(x_j)=\col(x_i)}} s_{ij}\,x_j ,
\]
y dejamos $x_j'=x_j$ para $j\neq i$. Verifiquemos tres cosas.

Todos los sumandos de $x_i'$ son
monocromáticos del mismo color que $x_i$, luego su suma también lo es (sus
coordenadas no nulas viven entre las coordenadas de ese color), y $x_i'\neq 0$ porque
su coeficiente en $x_i$ es $1$. Además el cambio $(x_j)\mapsto(x_j')$ es triangular
con unos en la diagonal, así que $(x_j')$ es una base.

La igualdad
$b_p=x_p+\sum_{j>p}s_{pj}x_j$ involucra únicamente vectores $x_j$ con $j\ge p>i$, y
para esos índices $x_j'=x_j$: la fórmula no cambia.

Despejando
\[
  x_i=x_i'-\sum_{j>i,\,\col(x_j)=\col(x_i)}s_{ij}x_j'
\]
y reemplazando en
\eqref{eq:tri},
\[
  b_i=\Bigl(x_i'-\!\!\sum_{\substack{j>i\\ \col(x_j)=\col(x_i)}}\!\! s_{ij}x_j'\Bigr)
      +\sum_{j>i}s_{ij}x_j'
   \;=\; x_i'+\sum_{\substack{j>i\\ \col(x_j)\neq\col(x_i)}} s_{ij}\,x_j' \;\in\; U_i .
\]

Los índices $p<i$ pueden haber cambiado, pero siguen cumpliendo una relación de la
forma $b_p=x_p'+\sum_{j>p}s_{pj}'x_j'$, es decir, se conserva \eqref{eq:tri}; y les
llegará su turno. Recorriendo $i=n,n-1,\dots,1$ obtenemos una base monocromática con
todos los índices buenos, que es exactamente lo pedido. 

\section*{Observaciones}

\begin{enumerate}
  \item La primera parte es fácil y, sobre todo, funciona como sugerencia:
    avisa que el orden de los $x_i$ es parte de lo que hay que construir, que
    es justamente la entrada al problema. Y además muchos participantes pueden
    tener algún punto. Si el problema no incluye este ítem, se vuelve bastante
    más duro. Para la versión difícil basta con quedarse solamente con el
    segundo ítem: no hay que cambiar nada. También puede dejarse solo el
    segundo ítem y poner como sugerencia que conviene empezar por reordenar la
    base canónica. 

  \item Los vectores $b_1,\dots,b_t$ construidos en el paso 2 son una base de
    $U_t$ para cada $t$ (son $t$ vectores independientes dentro de un espacio de
    dimensión $t$). Así, el segundo ítem equivale a lo siguiente: \emph{dada una
      bandera completa $U_1\subset\cdots\subset U_n$ y una descomposición
      $\C^n=V(1)\oplus\cdots\oplus V(l)$, existen una base $(x_i)$ con cada $x_i$
      en algún $V(k)$, una base $(b_i)$ con $U_t=\gen{b_1,\dots,b_t}$, y una
      matriz $S=(s_{ij})$ triangular superior con unos en la diagonal y con
      $s_{ij}=0$ toda vez que $x_i$ y $x_j$ estén en el mismo sumando, tales que
    $b_i=\sum_j s_{ij}x_j$.} Pintar las coordenadas no es un caso particular:
    eligiendo una base adaptada a la descomposición, todo $V(k)$ se convierte en
    un subespacio de coordenadas. Sin la condición sobre los colores, esto es la
    descomposición $LU$ de una matriz inversible después de permutar columnas.
\end{enumerate}


\end{document}
